a) Lingkaran besar adalah

\mathdisp {G= \left\{ (x,y,z) \in S  \mid  ax+by+cz = 0        \right\}} {  }

dan kedua hemisfer terbuka adalah

\mathdisp {U_+ = \left\{ (x,y,z) \in S  \mid  ax+by+cz > 0        \right\}} {  }

dan

\mathdisp {U_- = \left\{ (x,y,z) \in S  \mid  ax+by+cz < 0        \right\}} { . }

b) Untuk setiap hemisfer terbuka $U$ dan lingkaran besar $G$ yang terkait berlaku
\mathl{U \cap G = \emptyset}{.} Jarak maksimum antara dua titik
\mathl{P,Q \in S}{} adalah
\mathl{2}{,} dan hal ini terjadi tepat ketika kedua titik tersebut saling berhadapan
\zusatzklammer {antipodal} {} {}, yaitu ketika garis penghubungnya melalui pusat bola
\zusatzklammer {yakni pada
\mathl{Q=-P}{}} {} {}. Pasangan antipodal seperti itu tidak terletak pada satu hemisfer terbuka, sebab untuk
\mathl{P=(x,y,z)}{} dan
\mathl{P \in U_+}{} berlaku
\mathl{ax+by+cz>0}{} dan karena itu
\mathl{a(-x) +b(-y)+c(-z) <0}{} berlaku, sehingga
\mathl{-P \in U_-}{.}

Sekarang anggaplah bahwa terdapat suatu tutupan terbuka

\mathdisp {S \subseteq U_1 \cup U_2 \cup U_3} {  }

dengan tiga hemisfer terbuka
\mathl{U_1,U_2,U_3}{}
\zusatzklammer {bidang dan lingkaran besar yang bersesuaian masing-masing dilambangkan dengan
\mathkor {} {E_i} {atau} {G_i} {}} {} {.}
Karena
\mathl{U_1 \cap G_1 = \emptyset}{} diperoleh

\mathdisp {G_1 \subseteq U_2 \cup U_3} { . }

Irisan
\mathl{G_1 \cap E_2}{} memuat setidaknya dua titik antipodal
\mathkor {} {P} {dan} {-P} {}
dan
\mathl{P,-P \not\in U_2}{.} Karena, berdasarkan pengamatan sebelumnya, $U_3$ tidak memuat pasangan titik antipodal, salah satu dari kedua titik ini juga tidak termasuk dalam $U_3$, sehingga kita memperoleh kontradiksi.
